\begin{abstract}
Singapore's five-minute rainfall records offer a detailed view of local weather, while uneven coverage affects comparisons over time. We analyse the 2017--2024 annual datasets with Apache Spark, retaining 48.54 million observations from 91 station IDs after cleaning. Among sufficiently complete station-months, November has the highest mean monthly rainfall and wet intervals are most frequent in the afternoon. For prediction, we compare random forests, gradient-boosted trees, logistic regression and historical frequencies under two information settings: recent station rainfall and station/calendar information alone. On 5.73 million shared 2024 test cases, the recent-rainfall forest achieves mean Brier 0.01613, compared with 0.02283 for the calendar forest. Random forests give the lowest mean probability error among the three classifiers; historical frequency performs best for calendar-only risk. Increasing fitting samples from 15\% to 100\% produces small changes in the reported scores. Gradient boosting ranks 1 mm events competitively but overpredicts the average risk of large accumulations. The results favour recent-rainfall forests for short-term probability queries and identify high-threshold calibration as a useful direction for further modelling.
\end{abstract}

\section{Introduction}
The project asks two related questions: how rainfall varies across Singapore within and between years, and how well we can estimate a station's chance of receiving more than a specified amount in the next 30 minutes. The first question concerns patterns across many observations. The second requires a clear account of what information would be available when a forecast is issued.

We use all eight annual datasets from 2017 to 2024. For exploration, the main difficulty is uneven coverage: a low recorded monthly total can reflect missing observations as well as dry weather. We address this through completeness rules and comparisons at common stations. For prediction, we distinguish a calendar-only model from a model that also uses recently published rainfall. This lets us examine both the three-input query in the assignment---station, prediction time and threshold---and the benefit of supplying current weather observations internally.

The report follows these two questions from data preparation to results. Spark handles the observation-level calculations and model fitting; figures and supplementary analyses use the saved result tables.
